% Options for packages loaded elsewhere
\PassOptionsToPackage{unicode}{hyperref}
\PassOptionsToPackage{hyphens}{url}
%
\documentclass[
]{article}
\usepackage{amsmath,amssymb}
\usepackage{iftex}
\ifPDFTeX
  \usepackage[T1]{fontenc}
  \usepackage[utf8]{inputenc}
  \usepackage{textcomp} % provide euro and other symbols
\else % if luatex or xetex
  \usepackage{unicode-math} % this also loads fontspec
  \defaultfontfeatures{Scale=MatchLowercase}
  \defaultfontfeatures[\rmfamily]{Ligatures=TeX,Scale=1}
\fi
\usepackage{lmodern}
\ifPDFTeX\else
  % xetex/luatex font selection
\fi
% Use upquote if available, for straight quotes in verbatim environments
\IfFileExists{upquote.sty}{\usepackage{upquote}}{}
\IfFileExists{microtype.sty}{% use microtype if available
  \usepackage[]{microtype}
  \UseMicrotypeSet[protrusion]{basicmath} % disable protrusion for tt fonts
}{}
\makeatletter
\@ifundefined{KOMAClassName}{% if non-KOMA class
  \IfFileExists{parskip.sty}{%
    \usepackage{parskip}
  }{% else
    \setlength{\parindent}{0pt}
    \setlength{\parskip}{6pt plus 2pt minus 1pt}}
}{% if KOMA class
  \KOMAoptions{parskip=half}}
\makeatother
\usepackage{xcolor}
\usepackage{longtable,booktabs,array}
\usepackage{calc} % for calculating minipage widths
% Correct order of tables after \paragraph or \subparagraph
\usepackage{etoolbox}
\makeatletter
\patchcmd\longtable{\par}{\if@noskipsec\mbox{}\fi\par}{}{}
\makeatother
% Allow footnotes in longtable head/foot
\IfFileExists{footnotehyper.sty}{\usepackage{footnotehyper}}{\usepackage{footnote}}
\makesavenoteenv{longtable}
\usepackage{graphicx}
\makeatletter
\def\maxwidth{\ifdim\Gin@nat@width>\linewidth\linewidth\else\Gin@nat@width\fi}
\def\maxheight{\ifdim\Gin@nat@height>\textheight\textheight\else\Gin@nat@height\fi}
\makeatother
% Scale images if necessary, so that they will not overflow the page
% margins by default, and it is still possible to overwrite the defaults
% using explicit options in \includegraphics[width, height, ...]{}
\setkeys{Gin}{width=\maxwidth,height=\maxheight,keepaspectratio}
% Set default figure placement to htbp
\makeatletter
\def\fps@figure{htbp}
\makeatother
\ifLuaTeX
  \usepackage{luacolor}
  \usepackage[soul]{lua-ul}
\else
  \usepackage{soul}
\fi
\setlength{\emergencystretch}{3em} % prevent overfull lines
\providecommand{\tightlist}{%
  \setlength{\itemsep}{0pt}\setlength{\parskip}{0pt}}
\setcounter{secnumdepth}{-\maxdimen} % remove section numbering
\ifLuaTeX
  \usepackage{selnolig}  % disable illegal ligatures
\fi
\usepackage{bookmark}
\IfFileExists{xurl.sty}{\usepackage{xurl}}{} % add URL line breaks if available
\urlstyle{same}
\hypersetup{
  hidelinks,
  pdfcreator={LaTeX via pandoc}}

\author{}
\date{}

\begin{document}

\textbf{Multi-Agent Learning \& Imperfect Information Games:}

\textbf{Final Report Encompassing \emph{Spring 2025} Semester}

Mauricio Garcia

Supervised by Dr. Tian

University of Texas - Austin

April 2025

\section*{}\label{section}
\addcontentsline{toc}{section}{}

\subsection*{}\label{section-1}
\addcontentsline{toc}{subsection}{}

\subsection*{Introduction}\label{introduction}
\addcontentsline{toc}{subsection}{Introduction}

Imperfect‑information games, in which each participant holds private
information and must infer or bluff about hidden details, pose unique
challenges for AI. Traditional methods from perfect‑information domains
(for example, brute‑force search in chess) cannot directly handle hidden
information or deception. Instead, I adopted Counterfactual Regret
Minimization (CFR), an iterative self‑play algorithm that converges to
Nash equilibrium strategies for extensive‑form games. CFR has become
fundamental in solving poker and other bluffing games, automatically
revealing sophisticated tactics without explicit programming.

This report opens with a walkthrough of CFR on the classic Kuhn Poker
toy game. After experimenting with everything from Rock Paper Scissors
to Leduc and Kuhn Poker---and \emph{\ul{many}} of Liar's Dice
versions---I decided to focus this report on one custom three-stage
Liar's Dice game (bid → response → truth/lie reveal) as my
imperfect-information testbed. I describe the design of that game,
detail my CFR implementation (including state abstraction, action
pruning, opponent‑modeling tweaks, and a ``blueprint\,+\,subgame
resolving'' approach inspired by Sandholm's Libratus \& Pluribus
system), and evaluate the resulting strategies through both self‑play
and matchups against fixed‑strategy opponents.

Finally, I discovered that while CFR excels at finding unexploitable
equilibrium strategies, it does not natively deviate from equilibrium to
exploit human play patterns without extensive fine-tuning---a challenge
I wasn't able to overcome in my custom Liar's Dice variant, given the
difficulty and project's time constraints.

\section*{\texorpdfstring{\textbf{Background: CFR and Kuhn
Poker}}{Background: CFR and Kuhn Poker}}\label{background-cfr-and-kuhn-poker}
\addcontentsline{toc}{section}{\textbf{Background: CFR and Kuhn Poker}}

\textbf{Counterfactual Regret Minimization (CFR)} is an iterative
algorithm for finding approximate Nash equilibria in extensive-form
games with imperfect information. The key idea is to have the players
self-play the game repeatedly, adjusting their strategies based on
regret for suboptimal choices in hindsight. After each iteration, if a
player realizes that an alternative action would have yielded a better
outcome given the opponents' play, the algorithm increases the
propensity of that action (this is regret matching). An engine was made
for this, which I'll go into more detail about later.

Over many iterations, the strategy frequencies converge such that no
player has significant regret for any single deviation, meaning the
strategy profile approximates a Nash equilibrium. A crucial result in
CFR is that the average strategy converges to a Nash equilibrium as
cumulative regret diminishes. This allows strong, unexploitable
strategies to be learned even in games where explicit calculation of
equilibrium is intractable. For instance, heads-up no-limit Texas
Hold'em, the benchmark imperfect-information poker game, has on the
order of 10¹⁶⁵ distinct game states once you account for all possible
private cards and betting sequences. There's no hope of writing down and
solving the exact Nash equilibrium for a game that size by traditional
methods. Yet in 2017, Libratus---built on a CFR-based blueprint plus
subgame resolving approach---computed an approximate equilibrium
strategy for that same massive game and went on to defeat top human pros
by a statistically overwhelming margin. That result shows that CFR can
learn extremely strong, effectively unexploitable strategies, even when
explicit equilibrium computation is mathematically and computationally
intractable.

A toy example for CFR is \textbf{Kuhn Poker}, a simple two-player poker
game introduced by Harold Kuhn. Kuhn poker uses a deck of just 3 cards
(typically the Jack, Queen, King) where each player antes (small forced
bet to enter a Poker game) 1 chip and receives one private card. There
is a single round of betting: players can either pass (fold/check) or
bet (wager an extra chip). If one player bets and the other passes, the
bettor wins the pot immediately; if both players pass or both bet, they
proceed to a showdown where the higher card wins the pot.

Despite its small size (only 3 cards leading to 12 information sets in
total), Kuhn poker exhibits the essential features of poker: players
must bluff with weak hands and sometimes slow-play (i.e. check) strong
hands to maximize value. In fact, Kuhn poker was among the first games
analyzed in game theory and was solved analytically by hand in 1950. Its
equilibrium strategy is intuitive yet instructive; in equilibrium, when
you hold the worst card, you'll bluff (i.e. bet) only with a small
probability, and when you hold the best card, you'll bet almost every
time. Those exact bluff-and-bet probabilities are chosen so that
whenever your opponent faces a bet, their expected payoff from calling
equals their payoff from folding. CFR can be applied to Kuhn poker as a
didactic example (credit to Sandholm's Computational Game Solving course
at CMU), iteratively converging to this equilibrium strategy.

Implementing CFR on Kuhn poker demonstrates how the algorithm
automatically learns to bluff and call optimally: classic behaviors like
bluffing arise naturally from the regret minimization dynamics. This
small-scale success illustrates CFR's relevance to solving larger
imperfect-information games: if it can solve Kuhn's mini-poker, then
with enhancements it can tackle far more complex games like Texas
Hold'em. Indeed, modern poker AIs use CFR at their core -- the best
poker programs are now generated via self-play algorithms based on
game-theoretic principles, a leap from the early days when even large
poker games were approached with rule-based systems. CFR provides a
systematic way to approach optimal play, making it a natural choice for
designing an AI for any game of hidden information and strategic
bluffing.

\section*{\texorpdfstring{\textbf{Game Design: A Custom 3-Stage Liar's
Dice
Variant}}{Game Design: A Custom 3-Stage Liar's Dice Variant}}\label{game-design-a-custom-3-stage-liars-dice-variant}
\addcontentsline{toc}{section}{\textbf{Game Design: A Custom 3-Stage
Liar's Dice Variant}}

To explore CFR in a richer \& personalized setting, we designed a custom
game inspired by \textbf{Liar's Dice}, a classic dice bluffing game. In
Liar's Dice, multiple players roll dice in secret and take turns
claiming how many dice of a certain value appear under all cups, often
exaggerating (bluffing) their claims. Others can either accept the claim
or challenge it by calling the bluff (``liar!''), at which point, dice
are revealed to determine the truth. We created a simplified two-player,
single-round variant of Liar's Dice with three distinct stages:
\textbf{Bid}, \textbf{Response}, and \textbf{Truth/Lie Reveal}. This
3-stage structure captures the essence of one cycle of bluffing and
calling, while being self-contained enough for CFR analysis. The rules
and structure of the game are as follows:

\subsection{\texorpdfstring{\textbf{Private Dice Roll (Setup):}
}{Private Dice Roll (Setup): }}\label{private-dice-roll-setup}

\begin{itemize}
\item
  At the start, each player rolls a set of dice in secret. In our
  implementation, both Player 1 (P1) and Player 2 (P2) receive a private
  hand of dice (5 dice each with 6 faces). The results are known only to
  the respective player, modeling private information. This secret roll
  by ``chance'' sets up an information asymmetry: each player knows
  their own dice but not the opponent's.
\end{itemize}

\subsection{\texorpdfstring{\textbf{Stage 1 -- Bid (P1's Claim):}
}{Stage 1 -- Bid (P1's Claim): }}\label{stage-1-bid-p1s-claim}

\begin{itemize}
\item
  In the classic Liar's Dice variant, each bid specifies both a quantity
  Q and a face value F, as in ``There are at least 3 dice showing face
  5.'' This means ``across all hidden dice, I claim you'll find three
  5's.'' In our simplified version, we drop the face‐value detail and
  only care about the maximum frequency---that is, the count of
  whichever face appears most often. A bid now reads:

  \begin{itemize}
  \item
    ``Max-frequency ≥ Q''
  \end{itemize}
\end{itemize}

\subsection{When making a bid, P1 sees their own dice and chooses Q by
weighing:}\label{when-making-a-bid-p1-sees-their-own-dice-and-chooses-q-by-weighing}

\begin{itemize}
\item
  Plausibility: A higher Q is less likely to be true (more of a bluff).
\item
  Reward: If P2 doesn't call, an unchallenged higher Q wins outright.
\end{itemize}

\begin{itemize}
\item
  We also prune any impossible bids, like a Q greater than the total
  number of dice or Q that P1's hand plus the maximum the opponent could
  hold can't reach. This keeps the decision set both small and logically
  consistent.
\end{itemize}

\subsection{\texorpdfstring{\textbf{Stage 2 -- Response (P2's Challenge
or Accept):}
}{Stage 2 -- Response (P2's Challenge or Accept): }}\label{stage-2-response-p2s-challenge-or-accept}

\subsection{Once P1 announces ``Max-frequency ≥ Q,'' P2 has two
choices:}\label{once-p1-announces-max-frequency-q-p2-has-two-choices}

\begin{itemize}
\item
  \textbf{Raise:} Increase the bid to \textbf{Q′ \textgreater{} Q},
  forcing P1 back on the defensive.
\item
  \textbf{Call BS:} Challenge P1's claim immediately.
\end{itemize}

\begin{itemize}
\item
  If P2's hand already contains close to Q of some face, then
  ``Max-frequency ≥ Q'' is quite plausible, and P2 is likelier to raise
  (or even pass by raising minimally). If P2's hand makes Q seem out of
  reach, you can call a bluff. It can be summarized as follows: ``Given
  my dice, how likely is it that some face reaches Q across both
  hands?''
\end{itemize}

\subsection{\texorpdfstring{\textbf{Stage 3 -- Truth/Lie Reveal
(Resolution):}}{Stage 3 -- Truth/Lie Reveal (Resolution):}}\label{stage-3-truthlie-reveal-resolution}

\begin{itemize}
\item
  \textbf{If P2 called BS in Stage 2}, all dice are checked, and you
  check the actual maximum frequency vs. the last bid \textbf{Q}:

  \begin{itemize}
  \item
    \textbf{Actual ≥ Q} → P2 was wrong to call, so P2 loses a die.
  \item
    \textbf{Actual \textless{} Q} → P1 was caught bluffing, so P1 loses
    a die.
  \end{itemize}
\item
  \textbf{If P2 raised in Stage 2}, P1 now must \textbf{declare}
  ``truth'' or ``lie'' about their own claim, then everyone reveals
  their dice, and you compare actual vs. Q:

  \begin{itemize}
  \item
    P1 says ``truth''

    \begin{itemize}
    \item
      If actual ≥ Q, P2 loses a die.
    \item
      If actual \textless{} Q, P1 loses a die.
    \end{itemize}
  \item
    P1 says ``lie''

    \begin{itemize}
    \item
      If actual \textless{} Q, P2 loses a die.
    \item
      If actual ≥ Q, P1 loses a die.
    \end{itemize}
  \end{itemize}
\end{itemize}

This final stage ties together the bluff and the call/raise, rewarding
correct calls and catches, and penalizing mistaken challenges or failed
bluffs. This version keeps the game tree tractable for CFR analysis. I
was not particularly content with the third stage, as it does not feel
like a particularly `fair' or `balanced' ending. It is also not
particularly intuitive or easy to understand. In future work, I'd like
to build on a more intuitive game (like normal Liar's dice, etc.).

\section*{\texorpdfstring{\textbf{CFR Training with a Unified
Engine}}{CFR Training with a Unified Engine}}\label{cfr-training-with-a-unified-engine}
\addcontentsline{toc}{section}{\textbf{CFR Training with a Unified
Engine}}

\subsection*{\texorpdfstring{ The CFR logic is implemented once in
cfr\_train.py using a generic CFR class (from stub.py, obtained and
developed from Sandholm) parameterized by RegretMatchingPlus. Both
Player 1 and Player 2 policies are computed in the same loop. Here's the
high-level
workflow:}{ The CFR logic is implemented once in cfr\_train.py using a generic CFR class (from stub.py, obtained and developed from Sandholm) parameterized by RegretMatchingPlus. Both Player 1 and Player 2 policies are computed in the same loop. Here's the high-level workflow:}}\label{the-cfr-logic-is-implemented-once-in-cfr_train.py-using-a-generic-cfr-class-from-stub.py-obtained-and-developed-from-sandholm-parameterized-by-regretmatchingplus.-both-player-1-and-player-2-policies-are-computed-in-the-same-loop.-heres-the-high-level-workflow}
\addcontentsline{toc}{subsection}{ The CFR logic is implemented once in
cfr\_train.py using a generic CFR class (from stub.py, obtained and
developed from Sandholm) parameterized by RegretMatchingPlus. Both
Player 1 and Player 2 policies are computed in the same loop. Here's the
high-level workflow:}

\begin{enumerate}
\def\labelenumi{\arabic{enumi}.}
\item
  \textbf{Game Setup and Tuplification}

  \begin{itemize}
  \item
    Load the JSON‑encoded extensive-form game (decision nodes,
    observation mappings, and utility entries).
  \item
    Convert list-based edges and sequences into tuples for dictionary
    keys.
  \end{itemize}
\item
  \textbf{CFR+ Iterations}
\end{enumerate}

\subsection{Instantiate two CFR
solvers:}\label{instantiate-two-cfr-solvers}

\begin{enumerate}
\def\labelenumi{\arabic{enumi}.}
\item
  cfr1 = Cfr(dp1, rm\_class=RegretMatchingPlus)
\item
  cfr2 = Cfr(dp2, rm\_class=RegretMatchingPlus)

  \begin{enumerate}
  \def\labelenumii{\arabic{enumii}.}
  \item
    where dp1 and dp2 are the decision-problem lists for P1 (bid/truth)
    and P2 (respond).
  \end{enumerate}
\end{enumerate}

\begin{itemize}
\item
  Maintain cumulative sequence-form strategy sums (cum1, cum2) to build
  the average strategy over iterations.
\end{itemize}

\subsection{\texorpdfstring{\textbf{Loop} for t in 1..N iterations (for
example, 50k which I
used):}{Loop for t in 1..N iterations (for example, 50k which I used):}}\label{loop-for-t-in-1..n-iterations-for-example-50k-which-i-used}

\begin{enumerate}
\def\labelenumi{\arabic{enumi}.}
\item
  \textbf{Next Strategy}: Call cfr1.next\_strategy() and
  cfr2.next\_strategy(), yielding realization-probability vectors. These
  are mappings from each decision sequence to the probability of
  reaching that sequence under the current mixed strategy.
\item
  \textbf{Accumulate}: Add these (weighted by t) to cum1/cum2 to form
  the running average.
\item
  \textbf{Utility Computation}: Compute counterfactual utility vectors
  u1 and u2 with the opponent's current strategy.
\item
  \textbf{Regret Update}: Call cfr1.observe\_utility(u1) and
  cfr2.observe\_utility(u2) to update cumulative regrets at each node.
\end{enumerate}

\begin{enumerate}
\def\labelenumi{\arabic{enumi}.}
\setcounter{enumi}{2}
\item
  \textbf{Policy Extraction\\
  }After convergence, normalize the cumulative weights at each decision
  node to extract corresponding JSON policies.
\item
  \textbf{Unified Engine}
\end{enumerate}

\begin{quote}
A single CFR implementation serves both players: the same Cfr +
RegretMatchingPlus machinery handles P1's bid and truth/lie stages and
P2's response stage. This consistency avoids code duplication and
ensures identical update logic.
\end{quote}

\begin{enumerate}
\def\labelenumi{\arabic{enumi}.}
\setcounter{enumi}{4}
\item
  \textbf{Abstractions \& Pruning}

  \begin{itemize}
  \item
    \textbf{Action Pruning}
  \end{itemize}
\end{enumerate}

\subsection{\texorpdfstring{\textbf{State Abstraction}:
}{State Abstraction: }}\label{state-abstraction}

\begin{enumerate}
\def\labelenumi{\arabic{enumi}.}
\item
  Treat dice hands as multisets to collapse permutations into single
  info-sets.
\end{enumerate}

\subsection{\texorpdfstring{\textbf{Scaling Experiments}:
}{Scaling Experiments: }}\label{scaling-experiments}

\begin{enumerate}
\def\labelenumi{\arabic{enumi}.}
\setcounter{enumi}{1}
\item
  I trained on different factors, such as different starting dice
  amounts.
\end{enumerate}

\begin{enumerate}
\def\labelenumi{\arabic{enumi}.}
\setcounter{enumi}{5}
\item
  \textbf{Extensions \& Live Resolving}
\end{enumerate}

\subsection{\texorpdfstring{\textbf{Blueprint + Resolve}:
}{Blueprint + Resolve: }}\label{blueprint-resolve}

\begin{enumerate}
\def\labelenumi{\arabic{enumi}.}
\item
  Maintain the offline equilibrium (via ``blueprint'') and, in live
  play, analytically compute P(max ≥ Q \textbar{} P2's dice) for
  off-path bids.
\end{enumerate}

\subsection{\texorpdfstring{\textbf{Opponent Modeling}:
}{Opponent Modeling: }}\label{opponent-modeling}

\begin{enumerate}
\def\labelenumi{\arabic{enumi}.}
\setcounter{enumi}{1}
\item
  Freeze P2's strategy as ``always-call'' or ``always-accept'' (amongst
  others) to derive the best exploitative responses.
\end{enumerate}

The end result is three JSON policies mapping each info-set to an
equilibrium mixed strategy, ready for both large-scale self‑play
evaluation and interactive human-facing play.

stub.py provides a more general CFR engine (with both RegretMatching and
RegretMatchingPlus) used in coursework to explore theoretical
properties. It was from a homework assignment in Sandholm's
\href{https://www.cs.cmu.edu/~sandholm/cs15-888F23/}{\ul{course}} to
solve multiple problems given. cfr\_train.py was fully developed outside
of a course, and it uses RegretMatchingPlus CFR+ to compute full
equilibrium policies offline, focusing on convergence speed and
blueprint generation. It's more tailored for production, and not
necessarily classroom problems. For context, RegretMatching accumulates
all regrets (positive and negative), whereas RegretMatchingPlus (CFR+)
clamps any negative regret to zero each update, focusing only on
beneficial actions and speeding up convergence.

\section*{\texorpdfstring{\textbf{Evaluating Performance and
Scalability}}{Evaluating Performance and Scalability}}\label{evaluating-performance-and-scalability}
\addcontentsline{toc}{section}{\textbf{Evaluating Performance and
Scalability}}

\subsubsection*{\texorpdfstring{\textbf{Convergence to
Equilibrium}}{Convergence to Equilibrium}}\label{convergence-to-equilibrium}
\addcontentsline{toc}{subsubsection}{\textbf{Convergence to
Equilibrium}}

\subsection*{In the 3-stage version described, I managed to reach Nash
equilibrium:}\label{in-the-3-stage-version-described-i-managed-to-reach-nash-equilibrium}
\addcontentsline{toc}{subsection}{In the 3-stage version described, I
managed to reach Nash equilibrium:}

\begin{itemize}
\item
  Self-play win Rate ≈ 50 \%
\end{itemize}

After 50,000 CFR+ iterations on the full three-stage, 5-dice-per-player
variant, two independent AI instances split wins at 49--51 \% regardless
of who opened the bid. This symmetry shows neither side holds an
exploitable edge, i.e., the average policy approximates a Nash
equilibrium.

\subsubsection*{\texorpdfstring{\textbf{Exploiting an Always-Call
Opponent}}{Exploiting an Always-Call Opponent}}\label{exploiting-an-always-call-opponent}
\addcontentsline{toc}{subsubsection}{\textbf{Exploiting an Always-Call
Opponent}}

\begin{longtable}[]{@{}
  >{\raggedright\arraybackslash}p{(\columnwidth - 0\tabcolsep) * \real{1.0000}}@{}}
\toprule\noalign{}
\begin{minipage}[b]{\linewidth}\raggedright
=== 2-Stage Liar\textquotesingle s Dice Exploit Demo (3 dice each) ===\\
\strut \\
-\/-\/- Round 1 (AI: 3 dice, You: 3 dice) -\/-\/-\\
Your roll: {[}6, 4, 1{]}\\
AI opens ≥ 1\\
You: call\\
AI\textquotesingle s hand: {[}3, 2, 3{]} (true max-freq = 2)\\
→ You lose a die!\\
\strut \\
-\/-\/- Round 2 (AI: 3 dice, You: 2 dice) -\/-\/-\\
Your roll: {[}2, 5{]}\\
You open ≥ 2\\
AI: calls BS!\\
AI\textquotesingle s hand: {[}6, 4, 1{]} (true max-freq = 1)\\
→ You lose a die!\\
\strut \\
-\/-\/- Round 3 (AI: 3 dice, You: 1 die) -\/-\/-\\
Your roll: {[}4{]}\\
AI opens ≥ 1\\
You: call\\
AI\textquotesingle s hand: {[}6, 1, 5{]} (true max-freq = 1)\\
→ You lose a die!\\
\strut \\
*** AI wins the match! ***\strut
\end{minipage} \\
\midrule\noalign{}
\endhead
\bottomrule\noalign{}
\endlastfoot
\end{longtable}

\emph{This max-frequency style game follows the traditional two-round
until the highest bid approach. All bids ≥ 1 are doomed against a
caller; the AI finishes at a full 3-die advantage, playing a flawless
game.}

\subsubsection*{\texorpdfstring{\textbf{Vulnerability to Deviations\\
}Here we see that a seemingly innocuous deviation from the equilibrium
``always‐call'' benchmark---namely, always raising from Q=1 to
Q=2---causes the AI's blueprint to fail. In the five‐round ``deviation
demo,'' the AI loses every die in succession, wiping out its 3-die
starting advantage in a single
match:}{Vulnerability to Deviations Here we see that a seemingly innocuous deviation from the equilibrium ``always‐call'' benchmark---namely, always raising from Q=1 to Q=2---causes the AI's blueprint to fail. In the five‐round ``deviation demo,'' the AI loses every die in succession, wiping out its 3-die starting advantage in a single match:}}\label{vulnerability-to-deviations-here-we-see-that-a-seemingly-innocuous-deviation-from-the-equilibrium-alwayscall-benchmarknamely-always-raising-from-q1-to-q2causes-the-ais-blueprint-to-fail.-in-the-fiveround-deviation-demo-the-ai-loses-every-die-in-succession-wiping-out-its-3-die-starting-advantage-in-a-single-match}
\addcontentsline{toc}{subsubsection}{\textbf{Vulnerability to
Deviations\\
}Here we see that a seemingly innocuous deviation from the equilibrium
``always‐call'' benchmark---namely, always raising from Q=1 to
Q=2---causes the AI's blueprint to fail. In the five‐round ``deviation
demo,'' the AI loses every die in succession, wiping out its 3-die
starting advantage in a single match:}

\begin{longtable}[]{@{}
  >{\raggedright\arraybackslash}p{(\columnwidth - 0\tabcolsep) * \real{1.0000}}@{}}
\toprule\noalign{}
\begin{minipage}[b]{\linewidth}\raggedright
=== Deviation Demo: opponent always raises 1→2 instead of call ===

R1: AI={[}4, 1, 3{]} You={[}5, 5, 4{]} bid≥2 → You lose 1 die

R2: AI={[}3, 1, 1{]} You={[}4, 3{]} bid≥2 → You lose 1 die

R3: AI={[}1, 5, 3{]} You={[}4{]} bid≥2 → AI loses 1 die

R4: AI={[}6, 5{]} You={[}3{]} bid≥2 → AI loses 1 die

R5: AI={[}4{]} You={[}1{]} bid≥2 → AI loses 1 die

After deviation: AI has 0 dice, You have 1 dice.
\end{minipage} \\
\midrule\noalign{}
\endhead
\bottomrule\noalign{}
\endlastfoot
\end{longtable}

\emph{This reversal highlights a key tension in imperfect-information
AI: equilibrium strategies can be catastrophically fragile when the
opponent strays from the assumed distribution. In practice, then, pure
CFR must be augmented with opponent‐modeling or live subgame solving to
remain robust in the wild.}

\subsubsection*{\texorpdfstring{\textbf{Computational
Performance}}{Computational Performance}}\label{computational-performance}
\addcontentsline{toc}{subsubsection}{\textbf{Computational Performance}}

Using max-frequency (i.e. collapsing each hand to the single number
representing the highest count of any face) dramatically simplifies the
game's state space:

\subsection{Traditional face-specific Liar's
Dice:}\label{traditional-face-specific-liars-dice}

\begin{itemize}
\item
  Each player's 5-die hand is a histogram over faces 1--6 (e.g. 0 ones,
  2 twos, 1 three, 0 fours, 1 five, 1 six). There are
\item
  \href{https://www.codecogs.com/eqnedit.php?latex=\%5Cdisplaystyle\%20\%5Cbinom\%7B5\%20\%2B\%206\%20-\%201\%7D\%7B6\%20-\%201\%7D\%20\%3D\%20\%5Cbinom\%7B10\%7D\%7B5\%7D\%20\%3D\%20252\#0}{\includegraphics[width=1.83333in,height=0.375in]{media/media/image1.png}}
  possible frequency profiles (how many dice show each number).
\item
  \emph{5 is the number of dice you roll, 6 is the number of possible
  face values, and 1 comes from the standard
  combinations-with-repetition formula.}
\item
  Across both players, that's
  \href{https://www.codecogs.com/eqnedit.php?latex=\%5C(252\%5E2\%20\%5Capprox\%2063\%7B\%2C\%7D500\%5C)\#0}{\includegraphics[width=1.001in,height=0.15167in]{media/media/image2.png}}
  distinct chance-node outcomes before any bidding occurs
\end{itemize}

\subsection{Our custom max-frequency
adaptation:}\label{our-custom-max-frequency-adaptation}

\begin{itemize}
\item
  We reduce each hand to a single integer
  \href{https://www.codecogs.com/eqnedit.php?latex=\%5C(m\%20\%5Cin\%20\%5C\%7B1\%2C2\%2C3\%2C4\%2C5\%5C\%7D\%5C)\#0}{\includegraphics[width=1.23611in,height=0.15278in]{media/media/image3.png}}
  (the count of the most common face). That leaves just 5 possibilities
  per player, or
  \href{https://www.codecogs.com/eqnedit.php?latex=\%5C(5\%5E2\%20\%3D\%2025\%5C)\#0}{\includegraphics[width=0.59722in,height=0.16667in]{media/media/image4.png}}
  combined chance outcomes. This simple adaptation makes the game much
  more manageable and malleable.
\end{itemize}

In effect, max-frequency collapses the raw game tree by a factor of over
2,500× at the chance-node level alone, even before pruning bids or
responses. Downstream, the number of information sets across all three
stages shrinks from on the order of \textgreater100,000 (in a
face-specific model) to \textless≈12,000 with max-frequency abstraction.

\paragraph*{\texorpdfstring{\textbf{Practicality}}{Practicality}}\label{practicality}
\addcontentsline{toc}{paragraph}{\textbf{Practicality}}

From a practical standpoint, max-frequency abstraction yields
substantial efficiency gains without specialized hardware. By collapsing
each hand to a single integer, the memory footprint for storing regrets
and cumulative strategies shrinks from on the order of hundreds of
thousands of entries (≫ 1 GB) down to just a few thousand (≪ 1 GB). This
leaner state space also accelerates each CFR+ iteration by what seems
roughly 30--40 \%, which in turn reduces the total number of iterations
required to reach an average regret of ≤ 0.01 by a similar margin. On a
MacBook Pro M2 Pro with 16 GB of RAM, running 50,000 CFR+ iterations on
the full three-stage, 5-die max-frequency variant completes within a
reasonable window.

By abstracting, it is also easier to use interactive experimentation and
develop UIs, as you do not have to track the faces of the dice. However,
one of the trade-offs of this abstraction is clear: while we gain
dramatic tractability and a straightforward policy representation, we
forgo the ability to exploit opponents whose bluffing behavior hinges on
specific face distributions. Nevertheless, by preserving the core
dynamics of bluffing and calling, max-frequency abstraction strikes a
practical balance, enabling efficient CFR+ training on consumer-grade
hardware.

\section*{\texorpdfstring{\textbf{Key Lessons
Learned}}{Key Lessons Learned}}\label{key-lessons-learned}
\addcontentsline{toc}{section}{\textbf{Key Lessons Learned}}

\subsection{Equilibrium ≠ Exploitation (High
Win-rate):}\label{equilibrium-exploitation-high-win-rate}

\begin{itemize}
\item
  CFR+ converges to strategies with zero regret, but it does not
  automatically exploit suboptimal or non-equilibrium opponents. Without
  dedicated opponent modeling, the win rate will never be particularly
  high.
\end{itemize}

\subsection{Blueprint + Subgame
Resolving:}\label{blueprint-subgame-resolving}

\begin{itemize}
\item
  Storing an offline ``blueprint'' equilibrium and selectively resolving
  future decision points can adapt to opponent behavior and patch the
  brittleness seen in the deviation demo.
\end{itemize}

\subsection{Abstraction Trade-offs:}\label{abstraction-trade-offs}

\begin{itemize}
\item
  Collapsing dice permutations into multisets yields massive state‐space
  reductions, but over-abstraction risks throwing away exploitable
  patterns. Finding the right granularity is crucial.
\end{itemize}

\subsection{CFR+ Speeds Convergence:}\label{cfr-speeds-convergence}

\begin{itemize}
\item
  Clamping negative regrets to zero (using RegretMatchingPlus)
  accelerates convergence in large info‐sets, making CFR practical on a
  laptop for medium-sized games.
\end{itemize}

\section*{\texorpdfstring{\textbf{Conclusion}}{Conclusion}}\label{conclusion}
\addcontentsline{toc}{section}{\textbf{Conclusion}}

Over the course of this project, I've gained a deep appreciation for the
unique challenges posed by imperfect-information games and the power of
Counterfactual Regret Minimization to tackle them. I learned about a
problem I hadn't really thought through before---imperfect
information---and now I feel equipped to handle these scenarios a lot
better going forward. Designing and implementing a custom three-stage
Liar's Dice variant showed me how a simple max-frequency abstraction can
yield dramatic efficiency gains, while also highlighting the trade-offs
between tractability and exploitability. The biggest lesson was most
definitely that equilibrium strategies often stumble against real human
opponents unless you go the extra mile, like layering in live subgame
resolving or solid opponent modeling. It's not as simple as training a
Q-learning agent and walking away---you really have to lead it in the
right direction. Armed with these lessons, I'm excited for the future;
some work that can be done: refining my game design, deepening opponent
models, and deploying a lightweight UI so that I can test against human
players and iterate even faster.

\section*{\texorpdfstring{\textbf{References}}{References}}\label{references}
\addcontentsline{toc}{section}{\textbf{References}}

\begin{itemize}
\item
  Brown, N., \& Sandholm, T. (2018). Superhuman AI for heads-up
  non-limit poker: Libratus overcomes human talent. Science, 359(6374),
  418--424.
  \href{https://doi.org/10.1126/science.aao1733}{\ul{https://doi.org/10.1126/science.aao1733}}
\item
  Brown, N., \& Sandholm, T. (2019). Pluribus: The first AI to beat pros
  in heads-up, no-limit Texas hold 'em poker. Science, 365(6456),
  885--890.
  \href{https://doi.org/10.1126/science.aay2400}{\ul{https://doi.org/10.1126/science.aay2400}}
\item
  \href{https://www.youtube.com/watch?v=2dX0lwaQRX0&ab_channel=NoamBrown}{\ul{Superhuman
  AI for heads-up no-limit poker: Libratus beats top professionals}}
\item
  Kuhn, H. W. (1950). A simplified two-person poker. In H. W. Kuhn \& A.
  W. Tucker (Eds.), Contributions to the Theory of Games I (pp.
  97--103). Princeton University Press.
\item
  Neller, T. W., \& Lanctot, M. (2012). Counterfactual Regret
  Minimization: A tutorial. Technical Report, McGill University, Canada.
  Retrieved from
  \href{https://www.cs.toronto.edu/~vmnih/docs/draft-chapter.pdf}{\ul{https://www.cs.toronto.edu/\textasciitilde vmnih/docs/draft-chapter.pdf}}
\item
  Sandholm, T. (2021). Computational Game Theory and Equilibrium Finding
  (CMU CS 15-888 lecture notes). Carnegie Mellon University. Retrieved
  from
  \href{https://www.cs.cmu.edu/~sandholm/course/15-888-F21/lectures.html}{\ul{https://www.cs.cmu.edu/\textasciitilde sandholm/course/15-888-F21/lectures.html}}
\item
  \href{https://www.researchgate.net/publication/235737910_Measuring_the_Size_of_Large_No-Limit_Poker_Games}{\ul{https://www.researchgate.net/publication/235737910\_Measuring\_the\_Size\_of\_Large\_No-Limit\_Poker\_Games}}
\item
  \href{https://thegradient.pub/libratus-poker/}{\ul{https://thegradient.pub/libratus-poker/}}
\item
  \href{https://www.researchgate.net/publication/221605456_Effective_Short-Term_Opponent_Exploitation_in_Simplified_Poker}{\ul{https://www.researchgate.net/publication/221605456\_Effective\_Short-Term\_Opponent\_Exploitation\_in\_Simplified\_Poker}}
\end{itemize}

\end{document}
